\documentclass{article}
\usepackage[utf8]{inputenc}
\usepackage{amsmath}
\usepackage{amssymb}
\usepackage{geometry}
\geometry{a4paper, margin=1in}

\begin{document}

\begin{center}
  \textbf{SOAL ISIAN SINGKAT}
\end{center}

\begin{enumerate}
  \item Diberikan vektor tak nol $\mathbf{x} \in \mathbb{R}^{26}$. Misalkan $S = \{A \mid A \in M_{26 \times 26}(\mathbb{R}), A\mathbf{x} \in \text{span}(\mathbf{x})\}$. Nilai dari $\dim(S) = \dots$

  \item Bilangan dua digit $\overline{ab}$ dikatakan "jaya" jika matriks $A = \begin{bmatrix} a+5 & -b & 1 \\ b & a-5 & 2 \\ 0 & 0 & 3 \end{bmatrix}$ dapat didiagonalkan atas lapangan $\mathbb{R}$. Banyaknya bilangan "jaya" adalah $\dots$

  \item Nilai $\displaystyle \binom{2026}{2} + \binom{2026}{6} + \binom{2026}{10} + \dots + \binom{2026}{2026} = \dots$

  \item Diberikan $z_0 \in \mathbb{C} \setminus \{0\}$. Jika radius kekonvergenan deret $\displaystyle \sum_{n=0}^{\infty} c_n z^n$ adalah $R$ (dengan $0 < R < \infty$), maka radius kekonvergenan dari deret
        \[
          \sum_{n=0}^{\infty} (1 + z_0^n) c_n z^n
        \]
        adalah $\dots$

  \item Nilai dari $\displaystyle \lim_{n \to \infty} \frac{1}{n} \sqrt[n]{(n+1)(n+2)(n+3)\dots(n+n)}$ adalah $\dots$

  \item Jika fungsi $f : (0, 1] \to \mathbb{R}$ kontinu, tak konstan, dan memenuhi $f(x^2) = 2f(x)$ untuk setiap $x \in (0, 1]$, maka nilai dari
        \[
          \frac{1}{f\left(\frac{1}{2}\right)} \int_{\frac{1}{2}}^{1} \frac{f(x)}{x} dx
        \]
        adalah $\dots$

  \item Koefisien dari $x^{20}y^{26}$ pada ekspansi $\displaystyle \left(1 + \frac{1}{2}x + 2y\right)^{2026}$ adalah $\dots$

  \item Diberikan tiga himpunan $A$, $B$, dan $C$ dengan informasi sebagai berikut:
        \[
          |A| = |B| = |C| = 2026, \quad |A \cap B| = 202, \quad |A \cap C| = 206, \quad |B \cap C| = 226
        \]
        Nilai terbesar yang mungkin untuk $|A \cup B \cup C|$ adalah $\dots$

  \item Diberikan grup siklik $G$ dan $H$ dengan $|G| = 20$ dan $|H| = 26$. Banyak homomorfisma grup dari $G \times G$ ke $H$ adalah $\dots$

  \item Misalkan
        \[
          R = \left\{ \begin{pmatrix} a & b \\ 0 & c \end{pmatrix} : a,b,c \in \mathbb{Z}_6 \right\}
        \]
        adalah ring semua matriks segitiga atas berukuran $2 \times 2$ atas $\mathbb{Z}_6$. Banyak ideal dari $R$ adalah $\dots$
\end{enumerate}

\newpage

\begin{center}
  \textbf{SOAL URAIAN}
\end{center}

\vspace{0.5cm}

\begin{center}
  SUBBIDANG ALJABAR LINEAR
\end{center}

Diberikan transformasi linear $T : \mathbb{R}^n \to \mathbb{R}^n$. Buktikan bahwa terdapat bilangan positif $\alpha$ yang memenuhi
\[
  \|T(\mathbf{v})\| \le \alpha \|\mathbf{v}\|
\]
untuk setiap $\mathbf{v} \in \mathbb{R}^n$.

\vspace{0.5cm}

\begin{center}
  SUBBIDANG ANALISIS KOMPLEKS
\end{center}

Diberikan bilangan real $b_0, b_1, \dots, b_{2026}$ dengan $b_{2026} \neq 0$ dan $z_1, z_2, \dots, z_{2026}$ adalah akar-akar dari polinomial
\[
  P(z) = \sum_{k=0}^{2026} b_k z^k.
\]
Jika $\displaystyle \mu = \frac{|z_1| + \dots + |z_{2026}|}{2026}$ adalah rata-rata jarak dari $z_1, z_2, \dots, z_{2026}$ ke titik asal, tentukan konstanta terbesar $M$ sehingga $\mu \ge M$ untuk semua pilihan $b_0, b_1, \dots, b_{2026}$ yang memenuhi
\[
  1 \le b_0 < b_1 < b_2 < \dots < b_{2026} \le 2026.
\]

\vspace{0.5cm}

\begin{center}
  SUBBIDANG ANALISIS REAL
\end{center}

Carilah semua fungsi kontinu $f : \mathbb{R} \to \mathbb{R}$ dengan sifat terdapat fungsi $g : \mathbb{R}^2 \to \mathbb{R}$ yang memenuhi $f(xy) = g(x, f(y))$ untuk setiap $x, y \in \mathbb{R}$.

\vspace{0.5cm}

\begin{center}
  SUBBIDANG KOMBINATORIKA
\end{center}

Misalkan $a_1, a_2, \dots, a_{2026}$ adalah 2026 bilangan asli berbeda. Buktikan bahwa terdapat indeks $i, j$ dengan $1 \le i \le j \le 2026$ sedemikian sehingga jumlahan:
\[
  \sum_{k=i}^j (-1)^{k-i} a_k = a_i - a_{i+1} + a_{i+2} - \dots + (-1)^{j-i} a_j
\]
habis dibagi oleh $2026$.

\vspace{0.5cm}

\begin{center}
  SUBBIDANG STRUKTUR ALJABAR
\end{center}

Diberikan semigrup $(S, \star)$ yang memuat elemen identitas kiri $e$ yakni, $e \star x = x$ untuk setiap $x \in S$, dan berlaku sifat untuk setiap $y \in S$, terdapat $\bar{y} \in S$ dengan sifat $y \star \bar{y} = e$.
\begin{enumerate}
  \item[(a)] Berikan contoh di mana $(S, \star)$ belum tentu merupakan grup.
  \item[(b)] Jika elemen identitas kiri dari $(S, \star)$ tunggal, tunjukkan bahwa $(S, \star)$ merupakan grup.
\end{enumerate}

\end{document}